\section{Goal clarification before causal analysis and tool dispatch}
\label{exec23:gist-scope}
A request to improve a distressed-loan portfolio does not yet specify whether
the decision is to change contact frequency, offer a payment holiday, alter
collections or merely explain observed arrears. Dispatching a scientifically
valid estimator for the wrong intervention still answers the wrong business
question. GIST-CMTF addresses a related upstream problem: inferring a symbolic
goal before filtering the tools exposed to a language-model agent. Its June
2026 original develops candidate goals, ambiguity detection and clarification
within a common precondition--effect vocabulary. The authentic ten-page
original and a cutoff-selected repository revision are useful design evidence,
but their workflow use of causal does not by itself identify a treatment effect
or supply a structural causal model for borrowers.
\citep{exec23_gist_paper,exec23_gist_sources}

\subsection{What the state-transition model supplies}
\label{exec23:gist-model}
Represent a tool by $t_i=(d_i,R_i,E_i,c_i,\rho_i)$: a description, required
state predicates, produced predicates, optional cost and optional risk label.
A state $s$ and goal $g$ are sets of predicates. A symbolic tool is executable
when $R_i\subseteq s$, and its declared transition adds $E_i$ to the state.
The goal-inference layer maps the user's request and state to candidate goals
and confidence scores; the paper describes confidence and margin rules that
can request clarification instead of committing to one goal. Clarification is
itself represented as a transition, for example from an ambiguous request to
a specified goal. This supplies a concrete place to represent missing business
variables in a CLI or MCP request.
\citep{exec23_gist_paper}

The proposed banking extension should require an intervention, comparator,
population, outcome, horizon, information clock and decision owner before
constructing an effect query. A lender's goal predicate such as
\emph{servicing option selected} does not establish exchangeability, support
or a defensible adjustment set. Those are separate scientific obligations.
Likewise, a workflow edge saying a tool produces a graph does not establish
that the graph is causally correct. The useful inheritance is an explicit
contract for ambiguity and clarification, with scientific identification
review following goal selection.
\citep{exec23_gist_paper,exec23_gist_sources}

The inspected runner implements breadth-first search over monotone symbolic
states and returns the first tool on a shortest path, with ties determined by
tool order. Optional cost and risk labels in the conceptual representation
are not thereby optimized. The semantic-goal comparator uses lexical overlap
with candidate descriptions. The model inference prompt requests one selected
goal and its confidence, while the actual clarification rule combines that
output with a fixed confidence threshold, vague-word detection and the
benchmark's annotated request type. Thus the released implementation is a
specific hybrid rule, not an evaluation of an arbitrary calibrated goal
posterior or the complete candidate-margin formulation.
\citep{exec23_gist_sources}

\subsection{Reported results and reproducible aggregation}
\label{exec23:gist-results}
The paper evaluates 120 constructed tasks: 40 explicit goals, 40 ambiguous
goals, 30 missing-variable requests and 10 clarification-required requests.
Six methods and seven backends produce 5,040 saved rows, with deterministic
mocked tool effects. Independent data-only checks confirm 5,040 unique
model--method--task keys, 32 tools, positive saved token totals and the
published aggregate tables. GIST-CMTF succeeds in 815 of 840 rows, compared
with 673 for top-goal CMTF and 696 for the semantic-goal comparator. Its
21 wrong-goal outcomes contrast with 163 under top-goal CMTF. These are
reaggregations of the released final dataset, not independently repeated
model calls or evidence of general superiority on banking requests.
\citep{exec23_gist_paper,exec23_gist_sources,exec23_gist_reference}

The disaggregated results matter. On explicit-goal tasks the paper reports
94.6\% success for GIST-CMTF versus 98.9\% for the semantic-goal comparator,
with unnecessary clarification on 16.8\% of explicit requests. Its reported
advantage is concentrated in ambiguous requests. The final saved dataset
retains 170 labeled errors: 128 precondition failures, 24 tool-parse errors,
14 empty responses, three goal-parse errors and one invalid selection.
The authors state that infrastructure-corrupted rows were cleaned before the
final analysis. The final positive-token file cannot establish the complete
attempt, retry or provider-failure denominator without those original logs
and cleaning records. The paper explicitly leaves latency, monetary cost,
user satisfaction and real multi-turn clarification unmeasured.
\citep{exec23_gist_paper,exec23_gist_sources,exec23_gist_reference}

\subsection{A source-to-result reproducibility gap}
\label{exec23:gist-provenance}
The repository revision dated 30 June 2026 is later than the 15 June paper
but before the literature cutoff. Its identity is therefore distinct from
proof that it generated the published results. Eleven authenticated artifacts
include the runner, task generator, registry, task file, final metrics,
summary tables, citation metadata, dependency declaration and MIT notice.
The requirements are unpinned, and no raw provider trace appears in the
inspected complete tree. The paper's documented replacement of a Nova preview
backend with Nova Pro is visible in the saved rows, while the runner's default
list still names the preview profile. An explicit override could explain the
model list; it must be retained in a reproducible execution record.
\citep{exec23_gist_sources}

Three further discrepancies prevent direct replay claims. First, the runner
forces clarification for each nonexplicit request type and then assigns the
annotated intended goal. Yet 53 of 560 nonexplicit GIST rows in the saved data
report zero clarification. Second, the runner reads a \texttt{gold\_path}
field to count wrong tools. None of the 120 task objects has that field;
all instead provide \texttt{gold\_chains}. With those unmodified inputs,
that branch has no expected next tool against which to increment its count,
whereas the saved file contains 3,520 wrong-tool counts. Third, the runner's
written CSV columns differ from the released metrics columns. An unreleased
producer, transformation or earlier revision could reconcile these artifacts;
the available bytes do not establish which. This is an explicit provenance
gap, not an allegation that the reported observations were fabricated.
\citep{exec23_gist_sources,exec23_gist_reference}

Benchmark-assisted clarification also limits transport to real users. The
runner consults annotated request type and resolves a clarification by
substituting the intended goal, without observing a possibly ambiguous user
answer. The all-tools and state-aware comparators receive the intended goal
as prompt context. A production assessment must separately test learned
ambiguity detection, actual user responses and the absence of such labels at
inference time. Open-world goals, imperfect state predicates and failed or
nondeterministic tools remain explicit limitations in the original paper.
These limitations preserve the value of the design hypothesis while preventing
a claim that workflow goal validation solves general causal reasoning.
\citep{exec23_gist_paper,exec23_gist_sources}

\subsection{A public-evidence qualification route}
\label{exec23:gist-decision}
Retain GIST-CMTF as a goal-contract and evaluation reference. Before proposing
direct integration, reconcile the exact generating revision, task schema,
metric definitions, model overrides and cleaning log. A project-owned
prototype can then compare a label-free goal classifier, a clarification
policy, a simple lexical baseline and an explicitly named oracle on held-out
public or generated business requests. Freeze the scenario definitions before
outcome inspection; retain wrong goals, unsupported causal questions,
nonresponses, retries and abstentions in the full denominator. Actual
clarification responses and their costs must be represented, rather than
silently replacing them with the correct goal. This protocol is prospective;
no model is invoked or downloaded code executed for this source review.
\citep{exec23_gist_paper,exec23_gist_sources}

An illustrative cost equation makes the missing economics concrete. Let
$\Delta p$ be a descriptive difference in success frequency, $\Delta T$
the additional tokens and $\Delta C$ the additional clarifications per task.
For stipulated success value $v$, token price $c_T$ and clarification cost
$c_C$, define $\Delta N=v\Delta p-c_T\Delta T-c_C\Delta C$.
The saved GIST versus top-goal averages give
$\Delta p=71/420$, $\Delta T=45979/420$ and $\Delta C=23/70$.
With $v=10$ dollars and $c_T=1/1000$ dollar, clarification cost USD 2 gives
$\Delta N=388021/420000$ dollars, whereas cost USD 6 gives
$-163979/420000$ dollars. These exact stipulated amounts illustrate a decision
threshold; they do not estimate bank ROI or causal performance gains.
\citep{exec23_gist_reference}

The scientific owner remains responsible for translating a clarified business
objective into a causal estimand and diagram. The maintainer owns schema
compatibility, raw traces and versioned provenance; the operator owns the
actual endpoint, authentication and measured resource costs. A verified
workflow predicate should record what evidence established it. It must not
transform a language model's declaration that a goal or graph is correct into
an accepted scientific premise. The route is sufficiently concrete to justify
a bounded prototype proposal, while native capability, independent acceptance
and broader frontier coverage remain open.
